
\subsubsection{5.1 Pass Rates}

\begin{table}[htbp]
\centering
\resizebox{\columnwidth}{!}{%
\begin{tabular}{llll}
\toprule
Model & Pass & Fail & Pass Rate \\
\midrule
CodeRL & 306 & 236 & 56.5\% \\
CodeLlama-Instruct & 12 & 530 & 2.2\% \\
\bottomrule
\end{tabular}
}
\end{table}

The CodeRL model achieved substantially higher pass rates than CodeLlama-Instruct on these benchmarks. The low pass rate of CodeLlama-Instruct (2.2\%) indicates that this general instruction-following model performs poorly on these code generation tasks without additional code-specific fine-tuning.

\subsubsection{5.2 H-E1: Error Type Distribution Comparison}

\textbf{Contingency table (failure counts):}

\begin{table*}[htbp]
\centering
\resizebox{\dimexpr 2\columnwidth+\columnsep\relax}{!}{%
\begin{tabular}{lllll}
\toprule
 & Syntax & Runtime & Assertion & Total \\
\midrule
CodeRL & 218 & 12 & 5 & 236 \\
CodeLlama-Instruct & 529 & 1 & 0 & 530 \\
\bottomrule
\end{tabular}
}
\end{table*}

\textbf{Statistical results:}
\begin{itemize}
\item Chi-square = 35.27, df = 2
\item p = 2.19 $\times 10^{-8}$
\item Cramér's V = 0.215
\end{itemize}

The null hypothesis that error type distributions are independent of model is rejected at p < 0.001. The effect size (V = 0.215) is in the small-to-medium range by conventional standards.

\textbf{Proportions among failures:}

\begin{table}[htbp]
\centering
\resizebox{\columnwidth}{!}{%
\begin{tabular}{lll}
\toprule
Error Type & CodeRL & CodeLlama-Instruct \\
\midrule
Syntax & 92.4\% & 99.8\% \\
Runtime & 5.1\% & 0.2\% \\
Assertion & 2.1\% & 0.0\% \\
\bottomrule
\end{tabular}
}
\end{table}

\subsubsection{5.3 H-M1: Assertion Error Proportion}

\textbf{$2\times 2$ contingency table:}

\begin{table}[htbp]
\centering
\resizebox{\columnwidth}{!}{%
\begin{tabular}{llll}
\toprule
 & Assertion & Non-Assertion & Total \\
\midrule
CodeRL & 5 & 231 & 236 \\
CodeLlama-Instruct & 0 & 530 & 530 \\
\bottomrule
\end{tabular}
}
\end{table}

\textbf{Statistical results:}
\begin{itemize}
\item Fisher's exact test (one-sided, greater): p = 0.0027
\item CodeRL assertion proportion: 2.12\% (5/236)
\item CodeLlama-Instruct assertion proportion: 0.00\% (0/530)
\item Odds ratio: undefined (infinite, due to zero count)
\end{itemize}

The CodeRL model produced a non-zero proportion of assertion errors (2.12\%), while CodeLlama-Instruct produced zero assertion errors among 530 failures.

\subsubsection{5.4 H-M2: Execution Depth}

\textbf{Descriptive statistics:}

\begin{table*}[htbp]
\centering
\resizebox{\dimexpr 2\columnwidth+\columnsep\relax}{!}{%
\begin{tabular}{lllllll}
\toprule
Model & Mean Depth & SD & Median & Min & Max & N \\
\midrule
CodeRL & 0.294 & 0.311 & 0.211 & 0.0 & 1.0 & 236 \\
CodeLlama-Instruct & 0.001 & 0.022 & 0.0 & 0.0 & 0.5 & 530 \\
\bottomrule
\end{tabular}
}
\end{table*}

\textbf{Statistical results:}
\begin{itemize}
\item Welch's t-test: t = 14.47, p = 1.08 $\times 10^{-34}$
\item Cohen's d = 1.69
\end{itemize}

The CodeRL model's failures executed substantially deeper into the code (mean 29.4\% of lines) compared to CodeLlama-Instruct failures (mean 0.09\% of lines). The ratio of means is approximately 326:1. The effect size (d = 1.69) is large by conventional standards.

\textbf{95\% Confidence intervals:}
\begin{itemize}
\item CodeRL: [0.254, 0.334]
\item CodeLlama-Instruct: [-0.001, 0.003]
\end{itemize}

Trace success rate was 100\% for both models.

\subsubsection{5.5 H-M3: Fine-Grained Taxonomy}

\textbf{Coarse level (3 categories):}
\begin{itemize}
\item Chi-square = 33.79, df = 2
\item p = 4.61 $\times 10^{-8}$
\item Cramér's V = 0.210
\end{itemize}

\textbf{Fine-grained level (9 observed categories of 19 possible):}
\begin{itemize}
\item Chi-square = 525.40, df = 8
\item p = 2.49 $\times 10^{-108}$
\item Cramér's V = 0.823
\end{itemize}

The effect size increased from V = 0.21 at the coarse level to V = 0.82 at the fine-grained level.

\textbf{Fine-grained distribution:}

\begin{table}[htbp]
\centering
\resizebox{\columnwidth}{!}{%
\begin{tabular}{lll}
\toprule
Error Cause & CodeRL (n=236) & CodeLlama-Instruct (n=530) \\
\midrule
indentation\_error & 164 (69.5\%) & 0 (0.0\%) \\
syntax\_error & 54 (22.9\%) & 529 (99.8\%) \\
name\_error & 3 (1.3\%) & 1 (0.2\%) \\
type\_error & 5 (2.1\%) & 0 (0.0\%) \\
wrong\_output & 5 (2.1\%) & 0 (0.0\%) \\
value\_error & 2 (0.8\%) & 0 (0.0\%) \\
index\_error & 1 (0.4\%) & 0 (0.0\%) \\
recursion\_error & 1 (0.4\%) & 0 (0.0\%) \\
unknown & 1 (0.4\%) & 0 (0.0\%) \\
\bottomrule
\end{tabular}
}
\end{table}

CodeLlama-Instruct failures concentrated almost exclusively in a single error cause (syntax\_error: 99.8\%), while CodeRL failures distributed across multiple causes, with indentation\_error (69.5\%) and syntax\_error (22.9\%) most common.
